%% =====================================================================
%% PDF/A-2u (archival PDF, e.g. for the TU Delft repository) with the pdfx package.
%% Fonts embedded, sRGB output intent, all text mapped to Unicode, and XMP metadata.
%%
%% pdfx reads the metadata from a file named after the output (\jobname.xmpdata; on Overleaf the
%% output is called "output"). It is written here from _thesisinfo.tex on every compile, so there is
%% nothing to edit in it: change _thesisinfo.tex instead.
%%
%% Why not \DocumentMetadata: since 2025, \DocumentMetadata always loads LaTeX's new tagging code,
%% which replaces parts of memoir (captions, section headings, appendices). pdfx gives the same
%% PDF/A-2u without it, identically on every TeX Live version.
%% Needs: _thesisinfo.tex loaded first. Optional: \thesispdfsuffix (e.g. " (cover)") for the title.
%% =====================================================================
\chardef\thesisatcatcode=\catcode`\@ % restored at the end (this file is read inside the class too)
\makeatletter
\providecommand*{\thesispdfsuffix}{}
\providecommand*{\thesis@bcp}{en-GB}
\newwrite\thesis@xmpfile
\immediate\openout\thesis@xmpfile=\jobname.xmpdata
\immediate\write\thesis@xmpfile{\string\Title{\thesispdftitle\thesispdfsuffix}}
\immediate\write\thesis@xmpfile{\string\Author{\thesisauthorfirstnames\space\thesisauthorlastname}}
\immediate\write\thesis@xmpfile{\string\Subject{Doctoral dissertation, Delft University of Technology}}
\immediate\write\thesis@xmpfile{\string\Keywords{\thesiskeywords}}
\immediate\write\thesis@xmpfile{\string\Language{\thesis@bcp}}
\immediate\write\thesis@xmpfile{\string\Publisher{Delft University of Technology}}
\immediate\write\thesis@xmpfile{\string\Copyright{Copyright \the\year\space\thesisauthorfirstnames\space\thesisauthorlastname}}
\immediate\closeout\thesis@xmpfile
%% pdfx loads xcolor with an option ("hyperref") that newer xcolor versions declare obsolete and warn
%% about; remove just that option when xcolor loads.
\ExplSyntaxOn
\cs_new_protected:Npn \__thesis_drop_option:nn #1#2 % #1: option list macro name, #2: option
  {
    \cs_if_exist:cT {#1}
      {
        \clist_clear:N \l_tmpa_clist
        \exp_args:Nv \clist_map_inline:nn {#1}
          { \str_if_eq:eeF { \tl_trim_spaces:n {##1} } {#2} { \clist_put_right:Nn \l_tmpa_clist {##1} } }
        \cs_set_eq:cN {#1} \l_tmpa_clist
      }
  }
\AddToHook{package/xcolor/before}
  {
    \__thesis_drop_option:nn { opt@xcolor.sty } { hyperref }
    \__thesis_drop_option:nn { @raw@opt@xcolor.sty } { hyperref }
  }
\ExplSyntaxOff
%% PDF version 1.4 (allowed by PDF/A-2): hyperref's PDF/A mode, which pdfx switches on, expects it.
%% XeLaTeX (not the default engine) can't set the creation date the way PDF/A wants; pdfx then
%% warns instead of stopping.
\ifdefined\XeTeXrevision \PassOptionsToPackage{noerr}{pdfx}\fi
%% pdfx adds the font-map entry "dummy-space", which TeX Live's map already has; under LuaTeX it cannot hide
%% the resulting "duplicate entry" warning, so skip that one line while pdfx loads. (pdfx may load
%% luatex85, which defines \pdfmapline afresh: filter it again then.)
\ifdefined\directlua
    \def\thesis@filtermapline{\let\thesis@pdfmapline\pdfmapline
        \protected\def\pdfmapline##1{\in@{dummy-space}{##1}\ifin@\else\thesis@pdfmapline{##1}\fi}}
    \AddToHook{package/pdfx/before}{\thesis@filtermapline}
    \AddToHook{package/luatex85/after}{\thesis@filtermapline}
    \AddToHook{package/pdfx/after}{\let\pdfmapline\thesis@pdfmapline}
\fi


%% Newer pdfx reports "Setting all color commands to rgb" on every PDF/A compile. The template's colours
%% are all RGB or grey already, so that message is only information: show it as such in the log.
\let\thesis@PackageWarning\PackageWarning
\def\PackageWarning#1#2{%
    \in@{Setting all color commands}{#2}%
    \ifin@ \PackageInfo{#1}{#2}\else \thesis@PackageWarning{#1}{#2}\fi}
\let\thesis@PackageWarningNoLine\PackageWarningNoLine
\def\PackageWarningNoLine#1#2{%
    \in@{Setting all color commands}{#2}%
    \ifin@ \PackageInfo{#1}{#2}\else \thesis@PackageWarningNoLine{#1}{#2}\fi}
\RequirePackage[a-2u, pdf14]{pdfx}
\let\PackageWarning\thesis@PackageWarning
\let\PackageWarningNoLine\thesis@PackageWarningNoLine
\catcode`\@=\thesisatcatcode
